\documentclass[12pt]{article}
\usepackage[T2A]{fontenc}
\usepackage[utf8]{inputenc}
\usepackage[russian]{babel}
%\usepackage{amsfonts}
\usepackage{amssymb}
\usepackage{amsmath}
\usepackage{amsthm}
\usepackage{ccfonts,eulervm,microtype}
\renewcommand{\bfdefault}{sbc}

\usepackage[margin=0.7in,bmargin=1.2in]{geometry}
\usepackage{multicol}

\usepackage[colorlinks=false,pagebackref=true]{hyperref}

\usepackage{mathabx}

\usepackage{tikz-cd}
\usepackage{tikz}
\usetikzlibrary{arrows.meta,calc}

\pagestyle{plain}

\theoremstyle{plain}
\newtheorem{theorem}{Th.}[section]
\theoremstyle{lemma}
\newtheorem{lemma}{Лемма}
\theoremstyle{corollary}
\newtheorem{corollary}{Следствие}
\theoremstyle{proposition}
\newtheorem{proposition}{Предложение}
\theoremstyle{exercise}
\newtheorem{exercise}{Упражнение}


\theoremstyle{example}
\newtheorem{example}{Ex.}
\newtheorem{examples}[example]{Примеры}
\theoremstyle{remark}
\newtheorem{remark}{NB}

\theoremstyle{definition}
\newtheorem{definition}{Опр.}


\renewcommand{\emptyset}{\varnothing}
\newcommand\mbZ{\mathbb Z}
\newcommand\ph{\varphi}
\newcommand\trleq{\trianglelefteq}
\newcommand\isom{\cong}
%\def\l{\lambda}
%\def\m{\mu}
\newcommand\la{\langle}
\newcommand\ra{\rangle}
\newcommand\mb{\mathbb}
\newcommand\mc{\mathcal}
\newcommand\divs{\,\lower.4ex\vdots\,}
\newcommand\ol{\overline}
\newcommand\eps{\varepsilon}

\DeclareMathOperator{\ev}{ev}
\DeclareMathOperator{\id}{id}
\DeclareMathOperator{\Ker}{Ker}
\DeclareMathOperator{\Ree}{Re}
\DeclareMathOperator{\Img}{Im}
\DeclareMathOperator{\Arg}{Arg}
\DeclareMathOperator{\End}{End}
\DeclareMathOperator{\Aut}{Aut}
\DeclareMathOperator{\GL}{GL}
\DeclareMathOperator{\SL}{SL}
\DeclareMathOperator{\Hom}{Hom}
\DeclareMathOperator{\sgn}{sgn}
\DeclareMathOperator{\ord}{ord}
\DeclareMathOperator{\mmod}{mod}
\DeclareMathOperator{\cchar}{char}

\DeclareMathOperator{\logn}{ln}
\DeclareMathOperator{\Logn}{Ln}
\DeclareMathOperator{\Frac}{Frac}

\DeclareMathOperator{\inv}{inv}
\DeclareMathOperator{\adj}{adj}
\DeclareMathOperator{\rk}{rk}
\DeclareMathOperator{\pr}{pr}

\DeclareMathOperator{\pow}{pow}
%\DeclareMathOperator{\deg}{deg}
\DeclareMathOperator{\Fix}{Fix}

\DeclareMathOperator{\Map}{Map}
\DeclareMathOperator{\const}{const}


\newcommand\tld{\widetilde}
\newcommand\rsa{\rightsquigarrow}
\newcommand\mbC{\mathbb C}
\newcommand\mbR{\mathbb R}

\newcommand\literature[1]{{\small{\sc Литература}: #1}}

\newcommand\dfn[1]{{\bf #1}}

\makeindex

%\includeonly{multilinear}
\setcounter{section}{0}

\begin{document}

\section{Тензорное поле}

\subsection{Структура тензорного расслоения}

Пусть $\langle M, \tld{\mathcal{A}}_M\rangle$ - гладкое вещественное мноообразие размерности $\dim M = n$ 
и пусть $T_PM$ - его касаетльное пространство в точке $P \in M$.

\begin{definition}
    \textbf{Тензором над} $T_PM$ будем называть элемент тензорного произведения
    \begin{equation*}
        T^r_sV_P = (V^\ast_P)^{\otimes r} \otimes V^{\otimes s}_P, \quad V_P = T_PM.
    \end{equation*}  
\end{definition}

\begin{example}
    Если $\alpha^1, \ldots, \alpha^r \in V_P^\ast$ и $v_1, \ldots, v_s \in V_P$, тогда
    \begin{equation*}
        \alpha^1 \otimes \ldots \otimes \alpha^r \otimes v_1 \otimes \ldots \otimes v_s \in T^r_sV_P
    \end{equation*}
\end{example}

\begin{remark}
    Пусть далее $\langle N, \tld{\mathcal{A}}_N\rangle$ - другое гладкое многообразие 
    и $\sigma: M \to N$ - гладкое отображение между $M$ и $N$. Определим для некоторой 
    точки $P \in M$ отображение $T_P \sigma_s: T^0_sV_P \to T^0_sW_{\sigma(P)}$, где 
    $W_Q = T_QN$ следующим образом:
    \begin{equation*}
        \forall v_1 \ldots, v_s \in V_P, \quad 
        v_1 \otimes \ldots \otimes v_s \mapsto T_P\sigma (v_1) \otimes \ldots \otimes T_P \sigma (v_s)
    \end{equation*}
    Если $\sigma$ - диффеоморфизм, тогда можно определить отборажение 
    $T_P \sigma^r: (T^r_0 V_P)^\ast \to (T^r_0 W_{\sigma(P)})^\ast$ следующим образом:
    \begin{equation*}
        \forall \alpha^1, \ldots, \alpha^r \in V_P^\ast, \quad 
        \alpha^1 \otimes \ldots \otimes \alpha^r \mapsto [T_P\sigma^\ast]^{-1} (\alpha^1) \otimes \ldots \otimes [T_P\sigma^\ast]^{-1}(\alpha^r)
    \end{equation*}
    где $T_P\sigma$ - линейное отображение касательного пространства в точке $P \in M$:
    \begin{equation*}
        T_P\sigma: T_PM \to T_{\sigma(P)}N, \quad (T_P\sigma)^\ast: (T_{\sigma(P)}N)^\ast \to (T_PM)^\ast 
    \end{equation*}
\end{remark}

\begin{remark}
    Если $(U,\varphi) \in \tld{\mathcal{A}}_M$ - карта на $M$, 
    тогда $(TU, T\varphi)$ - карта на касательном расслоении $TM$, а 
    $(T^\ast U, T^\ast \varphi)$ - карта на кокасательном расслоении $T^\ast M$
    \begin{gather*}
        TU = \bigcup_{P\in U} T_PM, \quad 
        T\varphi: TU \to T\mb R^n, \\
        T^\ast U = \bigcup_{P\in U} T^\ast_PM, \quad
        T^\ast \varphi: T^\ast U \to (\mb R^n)^\ast ,\quad 
        T^\ast \varphi = (T \varphi^\ast)^{-1}.
    \end{gather*}
\end{remark}

\begin{lemma}
    Расслоения $TM$ и $T^\ast M$ - гладкие многобразия, причем 
    \begin{equation*}
        \dim TM = \dim T^\ast M = 2 \cdot \dim M.
    \end{equation*}
\end{lemma}

\begin{definition}
    \textbf{Тензорным расслоением} типа $(r,s)$ называется множество 
    \begin{equation*}
        T^r_sM = \bigcup_{P\in M} T^r_s V_P.
    \end{equation*}
\end{definition}

\begin{lemma}
    Если $(U,\varphi) \in \tld{\mathcal{A}}_M$ - карта на $M$, тогда 
    $(T^r_s U, T\sigma^r_s)$ - карта на $T^r_sM$ и 
    \begin{gather*}
        T^r_sU = \bigcup_{P \in U} T^r_s V_P, \quad 
        T\sigma^r_s: T^r_sU \to T^r_s\mb R^n \\ 
        T\sigma^r_s: \alpha^1 \otimes \ldots \otimes \alpha^r \otimes v_1 \otimes \ldots \otimes v_s \mapsto 
        [T_P\sigma^\ast]^{-1} (\alpha^1) \otimes \ldots \otimes [T_P\sigma^\ast]^{-1} (\alpha^r) \otimes 
        T_P\sigma (v_1) \otimes \ldots \otimes T_P\sigma (v_s).
    \end{gather*}
\end{lemma}

\begin{theorem}
    Тензорное расслоение $T^r_s M$ - гладкое $n^{r+s}$-мерное многообразие.
\end{theorem}

\subsection{Координаты в тензорном расслоении}

Пусть $(U,\varphi) \in \tld{\mathcal{A}}_M$ - координатная карта на $M$, 
причем $\{x^i\}_{i=1}^n$ - координатные функции на $U$:
\begin{equation*}
    \forall P \in U \quad \varphi(P) = (x_P^1, x_P^2, \ldots, x_P^n).
\end{equation*}

\begin{remark}
    На касательно расслоении $TM$ отображение $\varphi = (x^1, x^2, \ldots, x^n)$
    индуцирует координатное отображение $T\varphi$, которое для каждого 
    $v_P \in TM$ возвращает набор: 
    \begin{equation*}
        [T\varphi](v_P) = [D_P \varphi](v_P) = (x^1_P, x^2_P, \ldots, x^n_P; \xi^1, \xi^2, \ldots, \xi^n), \quad 
        \xi^i = dx_P^i(v_P), \quad 
        v_P = \sum_{i=1}^n \xi^i \partial_{x^i}.
    \end{equation*}
    Базис $C^\infty(M)$-модуля $\chi(M)$ векторных полей на $M$ образует набор 
    $\{\partial_i = \partial_{x^i}\}_{i=1}^n$.
\end{remark}

\begin{remark}
    Аналогично, в кокасательном расслоении $T^\ast M$ индуцируется коордиантное отображение 
    $T^\ast \varphi$, которое для каждого $\alpha_P \in T^\ast M$ возвращает набор
    \begin{equation*}
        [T^\ast\varphi](\alpha_P) = (x^1_P, x^2_P, \ldots, x^n_P; \eta_1, \eta_2, \ldots, \eta_n), \quad 
        \eta_i = \widehat{\partial}_{x^i}(\alpha) = \alpha(\partial_{x^i}), \quad 
        \alpha = \sum_{i=1}^{n} \eta_i dx^i.
    \end{equation*}
    Базис $C^\infty(M)$-модуля $\chi^\ast(M)$ дифференциальных 1-форм образует набор $\{dx^i\}_{i=1}^n$.
\end{remark}

\begin{remark}
    Отобрежения $T\varphi$ и $T\varphi^\ast$ индуцируют, в свою очередь, координатное 
    отображение $T\varphi^r_s$ на $T^r_sM$, а именно: для каждого $\tau_P \in T^r_sM$
    будем иметь 
    \begin{gather*}
        \tau_P = \alpha^1 \otimes \ldots \otimes \alpha^r \otimes v_1 \otimes \ldots \otimes v_s, \\
        [T^r_s\varphi](\tau_P) = (x^1_P, x^2_P, \ldots, x^n_P; \{\eta^1_{i_1}\ldots\eta^r_{i_r}\xi^{j_1}_1 \ldots \xi^{j_s}_s\}).
    \end{gather*}
\end{remark}

\subsection{Тензорное поле}

\begin{definition}
    \textbf{Естестенной проекцией} тензороного расслоения $T^r_s(M)$ называется отображение 
    \begin{equation*}
        \pi: T^r_sM \to M, \quad \tau_P \mapsto P
    \end{equation*}
\end{definition}

\begin{definition}
    \textbf{Тензорным полем} называется гладкое сечение $\tau: M \to T^r_sM$ тензорного расслоения:
    \begin{equation*}
        \pi \circ \tau = id_M.
    \end{equation*}
\end{definition}

\begin{remark}
    В координатной карте $(U, \varphi) \in \tld{\mathcal{A}}_M$ и индуцированной ею карте 
    $(T^r_sU, T^r_s \varphi)$ на $T^r_sM$ представитель отображения $\tau$ имеет вид:
    \begin{equation*}
        \widehat{\tau} = T^r_s\varphi \circ \tau \circ \varphi^{-1}, \quad 
        \widehat{\tau}(x^1, \ldots, x^n) = t_{i_1, i_2, \ldots, i_r}^{j_1, j_2, \ldots, j_s}(x^1, \ldots, x^n) \cdot
        x^{i_1} \otimes \ldots \otimes dx^{i_r} \otimes \partial_{j_1} \otimes \ldots \otimes \partial_{j_s}
    \end{equation*} 
\end{remark}

\begin{remark}
    Множество тензорных полей типа $(r,s)$ на $M$ договоримся обозначать $\chi^r_s(M)$.
\end{remark}

\begin{remark}
    Базис $C^\infty(M)$-модуля $\chi^r_s(M)$ образует следующий набор тензорных полей 
    \begin{equation*}
        \{dx^{i_1} \otimes \ldots \otimes dx^{i_r} \otimes \partial_{j_1} \otimes \ldots \otimes \partial_{j_s}\}.
    \end{equation*}
\end{remark}

\begin{lemma}
    При преобразовании координат $\{x^i\} \to \{y^j\}$, индуцированном переходом от карты 
    $(U,\varphi)$ к карте $(V,\psi)$ набор функций $t_{i_1, i_2, \ldots, i_r}^{j_1, j_2, \ldots, j_s}$ 
    преобразуется следующим образом:
    \begin{equation*}
        \tld{t}_{k_1, k_2, \ldots, k_r}^{m_1, m_2, \ldots, m_s} = 
        t_{i_1, i_2, \ldots, i_r}^{j_1, j_2, \ldots, j_s} \mu^{i_1}_{k_1} \ldots \mu^{i_r}_{k_r}\nu_{j_1}^{m_1} \ldots \nu_{j_s}^{m_s}, \quad 
        \mu = D(\psi \circ \varphi^{-1}) = \left(\frac{\partial y^j}{\partial x^i}\right), \quad 
        \nu = \mu^{-1}.
    \end{equation*}
\end{lemma}

\subsection{Производная тензорного поля}

\begin{definition}
    \textbf{Тензорной производной} называется семейство отображений $D^r_s: \chi^r_s(M) \to \chi^r_s(M)$, 
    (которые для удобства договоримся обозначать просто $D$), удовлетворяющее следующим аксиомам:
    \begin{enumerate}
        \item $D(a \tau + b \omega) = a \cdot D(\tau) + b \cdot D(\omega)$ - линейность;
        \item $D \circ \mathcal{C} = \mathcal{C} \circ D$ - коммутирование со сверткой;
        \item $D(\tau \otimes \omega) = D(\tau) \otimes \omega + \tau \otimes D(\omega)$ - правило Лейбница.
    \end{enumerate}
\end{definition}

\begin{example}
    Пример производной - производная Ли вдоль векторного поля $v \in \chi(M)$:
    \begin{enumerate}
        \item $f \in C^\infty(M), \quad Df = \mathcal{L}_v(f) = v(f)$;
        \item $u \in \chi(M), \quad \mathcal{L}_v(u) = [v,u]$.
    \end{enumerate}
\end{example}

\begin{example}
    Пользуясь определением, найдем выражение для $D(\alpha)$, где $\alpha \in \chi^\ast(M)$. Пусть 
    $u \in \chi(M)$, тогда
    \begin{equation*}
        D[\alpha(u)] = (D \circ \mathcal{C})(\alpha \otimes u) = 
        \mathcal{C}(D(\alpha) \otimes u + \alpha \otimes D(u)) = 
        [D(\alpha)](u) + \alpha[D(u)].
    \end{equation*}
    Таким образом, окончательно получим:
    \begin{equation*}
        [D(\alpha)](u) = D[\alpha(u)] - \alpha[D(u)] = (D \circ \alpha - \alpha \circ D)(u).
    \end{equation*}
\end{example}

\begin{remark}
    Пусть $v \in \chi(M)$ и $D = \mathcal{L}_v$, тогда 
    \begin{equation*}
        [\mathcal{L}_v(\alpha)](u) = v[\alpha(u)] - \alpha([v,u]).
    \end{equation*}
\end{remark}

\begin{lemma}
    Производная Ли может быть обобщена на полилинейные отображения.
\end{lemma}

\begin{proof}
    Рассмотрим форму $\omega: X^{\times r} \times (X^\ast)^{\times s} \to \mb R$. Имеем:
    \begin{equation*}
        \omega(\alpha^1, \ldots, \alpha^r, u_1, \ldots, u_s) = 
        \mathcal{C}(\omega \otimes \alpha^1 \otimes \ldots  \otimes \alpha^r \otimes u_1 \otimes \ldots \otimes u_s).
    \end{equation*}
    Применив к обеим частям оператор $D$, получим:
    \begin{gather*}
        D[\omega(\ldots)] = [D\omega](\ldots) + \sum_{k=1}^{r} \omega[\ldots D(\alpha^k) \ldots] + 
        \sum_{m=1}^{s} \omega[\ldots D(u_m) \ldots].
    \end{gather*}
    Переходя обратно $D = \mathcal{L}_v$, окончательно получим:
    \begin{equation*}
        [\mathcal{L}_v(\omega)](\ldots) = \mathcal{L}_v[\omega(\ldots)] - 
        \sum_{k=1}^{r} \omega[\ldots \mathcal{L}_v(\alpha^k) \ldots] - 
        \sum_{m=1}^{s} \omega[\ldots \mathcal{L}_v(u_m) \ldots].
    \end{equation*}
\end{proof}

\end{document}
